%------------------------------------------------------------------------------%
\begin{question}
  Utilize o escalonamento para resolver o sistema linear
  \[
    \systeme{
       x + 2y -  z = 3,
      2x + 5y +  z = 8,
      -x - y  + 2z = 1
      }
  \]
\end{question}

%------------------------------------------------------------------------------%
\begin{resolution}{Matriz Aumentada}
\[
  A = \left[
  \begin{array}{rrr}
    1  & 2  & -1  \\
    2  & 5  & 1   \\
    -1 & -1 & 2
  \end{array}
  \right]
  \pause
  \qquad\qquad
  b = \left[
  \begin{array}{rrr|r}
     3 \\
     8 \\
     1
  \end{array}
  \right]
\]

\pause
Matriz aumentada
\[
  \left[
  \begin{array}{rrr|r}
    1  & 2  & -1 & 3 \\
    2  & 5  & 1  & 8 \\
    -1 & -1 & 2  & 1
  \end{array}
  \right]
\]
\end{resolution}

%------------------------------------------------------------------------------%
\begin{resolution}{Escalonamento do Sistema}
\begin{columns}
\column{0.4\linewidth}
  \[\;\;\,
    \left[
    \begin{array}{rrr|r}
      1  & 2  & -1 & 3 \\
      2  & 5  & 1  & 8 \\
      -1 & -1 & 2  & 1
    \end{array}
    \right]
  \]
  \pause
\column{0.6\linewidth}
  Eliminar os elementos abaixo do primeiro pivô
  \\ \pause
  \(
    L_2 \leftarrow L_2 - 2 L_1
  \)
  \\ \pause
  \(
    L_3 \leftarrow L_3 + L_1
  \)
\end{columns}
\vspace{-\baselineskip}
\[
  \pause
  L'_2
  \;=\;
  L_2 - 2 L_1
  \pause
  \;=\;
  \left[
  \begin{array}{rrr|r}
    2  & 5  & 1  & 8
  \end{array}
  \right]
  -2
  \left[
  \begin{array}{rrr|r}
    1  & 2  & -1 & 3
  \end{array}
  \right]
  \pause
  \;=\;
  \left[
  \begin{array}{rrr|r}
    0 & 1 & 3  & 2
  \end{array}
  \right]
\]
\[
  \pause
  L'_3
  \;=\;
  L_3 + L_1
  \pause
  \;=\;
  \left[
  \begin{array}{rrr|r}
    -1 & -1 & 2  & 1
  \end{array}
  \right]
  +
  \left[
  \begin{array}{rrr|r}
    1  & 2  & -1 & 3
  \end{array}
  \right]
  \pause
  \;=\;
  \left[
  \begin{array}{rrr|r}
    0 & 1 & 1  & 4
  \end{array}
  \right]
\]

\pause
\[
  \left[
  \begin{array}{rrr|r}
    1 & 2 & -1 & 3 \\
    0 & 1 & 3  & 2 \\
    0 & 1 & 1  & 4
  \end{array}
  \right]
\]
\end{resolution}

%------------------------------------------------------------------------------%
\begin{resolution}{Escalonamento do Sistema}
\begin{columns}
\column{0.4\linewidth}
  \[\;\;\,
    \left[
    \begin{array}{rrr|r}
      1  & 2  & -1 & 3 \\
      2  & 5  & 1  & 8 \\
      -1 & -1 & 2  & 1
    \end{array}
    \right]
  \]
  \pause
\column{0.6\linewidth}
  Eliminar os elementos abaixo do segundo pivô
  \\ \pause
  \(
    L_3 \leftarrow L_3 - L_2
  \)
\end{columns}
\vspace{-\baselineskip}
\[
  \pause
  L'_3
  \;=\;
  L_3 - L_2
  \pause
  \;=\;
  \left[
  \begin{array}{rrr|r}
    -1 & -1 & 2  & 1
  \end{array}
  \right]
  -
  \left[
  \begin{array}{rrr|r}
    2  & 5  & 1  & 8
  \end{array}
  \right]
  \pause
  \;=\;
  \left[
  \begin{array}{rrr|r}
    0 & 0 & -2 & 2
  \end{array}
  \right]
\]

\pause
\[
  \left[
  \begin{array}{rrr|r}
    1 & 2 & -1 & 3 \\
    0 & 1 & 3  & 2 \\
    0 & 0 & -2 & 2
  \end{array}
  \right]
\]
\end{resolution}

%------------------------------------------------------------------------------%
\begin{resolution}{Sistema Linear Escalonado}
Matriz aumentada do sistema escalonado
\[
  \left[
  \begin{array}{rrr|r}
    1  & 2  & -1 & 3 \\
    \z & 1  & 3  & 2 \\
    \z & \z & -2 & 2
  \end{array}
  \right]
\]

\pause
Sistema linear escalonado
\[
  \systeme{
    x + 2y  -z = 3,
         y +3z = 2,
           -2z = 2}
\]
\end{resolution}

%------------------------------------------------------------------------------%
\begin{resolution}{Substituição Retroativa}
\begin{columns}[t]
\column{0.26\linewidth}
\onslide<+->{\centerline{Última linha}}
\begin{align*}
  \onslide<+->{-2z & = 2}  \\
  \onslide<+->{z   & = -1}
\end{align*}
\column{0.27\linewidth}
\onslide<+->{\centerline{Segunda linha}}
\begin{align*}
  \onslide<+->{y + 3z    & = 2} \\
  \onslide<+->{y + 3(-1) & = 2} \\
  \onslide<+->{y - 3     & = 2} \\
  \onslide<+->{y         & = 5}
\end{align*}
\column{0.4\linewidth}
\onslide<+->{\centerline{Primeira linha}}
\begin{align*}
  \onslide<+->{x + 2y - z          & = 3}  \\
  \onslide<+->{x + 2\times 5 -(-1) & = 3}  \\
  \onslide<+->{x + 10 + 1          & = 3}  \\
  \onslide<+->{x                   & = -8}
\end{align*}
\end{columns}
\end{resolution}

%------------------------------------------------------------------------------%
\begin{resolution}{Solução}
A solução do sistema é
\(
  {\bm x}
  \;=\;
  \begin{bmatrix}
    x \\ y \\ z
  \end{bmatrix}
  \;=\;
  \begin{bmatrix}
    -8 \\ 5 \\ -1
  \end{bmatrix}
\)
\end{resolution}

%------------------------------------------------------------------------------%
